\aufzaehlungvier{Dies ergibt sich aus der expliziten Beschreibung der vertikalen Ableitung zu einem linearen Zusammenhang, siehe
[[Offene Menge/Triviales Vektorbündel/Linearer Zusammenhang/Spaltung/Christoffelsymbole/Bemerkung|Bemerkung]].
}{Klar wegen (1).
}{Folgt aus (2) und aus
[[Differenzierbare Mannigfaltigkeit/R/Vektorbündel/Linearer Zusammenhang/Lokale Beschreibung/Fakt|Fakt]].
}{Wir schreiben

\mavergleichskette
{\vergleichskette
{ V
}
{ = }{ \sum_{i = 1}^n {h_i } \partial_i
}
{  }{
}
{    }{
}
{   }{
}
}
{}{}{}
und

\mavergleichskette
{\vergleichskette
{ W
}
{ = }{ \sum_{j = 1}^n {f_j } \partial_j
}
{  }{
}
{    }{
}
{   }{
}
}
{}{}{.}
Nach
[[Differenzierbare Mannigfaltigkeit/R/Vektorbündel/Linearer Zusammenhang/Lokale Beschreibung/Fakt|Fakt]]
und den schon bewiesenen Teilen ist

\mavergleichskettealign
{\vergleichskettealign
{ \nabla_V W
}
{ =} { \nabla_{\sum_{i = 1}^n {h_i } \partial_i }   \left(  \sum_{j = 1}^n {f_j } \partial_j   \right)
}
{ =} { \sum_{j = 1}^n \left(  \sum_{i = 1}^n h_i  \partial_i f_j   \right) \partial_j
}
{ } {
}
{ } {
}
}
{}
{}{.}
Für einen Punkt

\mavergleichskette
{\vergleichskette
{ P
}
{ \in }{ U
}
{  }{
}
{    }{
}
{   }{
}
}
{}{}{}
ist dies

\mathdisp {\sum_{j  = 1}^n \left(  \sum_{i = 1}^n h_i (P) ( \partial_i f_j)(P)   \right) \partial_j} { . }

Ebenso ist

\mavergleichskettealign
{\vergleichskettealign
{ \left(  D W   \right)_{ P } \left(   V(P)  \right)
}
{ =} { \left(  (\partial_j f_i)(P)  \right)_{ji} \begin{pmatrix} h_1(P) \\ \vdots\\  h_n(P)   \end{pmatrix}
}
{ =} { \begin{pmatrix}  \sum_{i = 1}^n h_i (P) \left(  \partial_i f_1  \right)(P)  \\ \vdots\\  \sum_{i = 1}^n h_i (P)  \left(  \partial_i f_n  \right) (P)   \end{pmatrix}
}
{ } {
}
{ } {
}
}
{}
{}{.}
}

 [[Kategorie:Latexseite]]
